% Options for packages loaded elsewhere
\PassOptionsToPackage{unicode}{hyperref}
\PassOptionsToPackage{hyphens}{url}
\PassOptionsToPackage{dvipsnames,svgnames,x11names}{xcolor}
\documentclass[
  11pt,
]{article}
\usepackage{xcolor}
\usepackage[margin=2.3cm]{geometry}
\usepackage{amsmath,amssymb}
\setcounter{secnumdepth}{-\maxdimen} % remove section numbering
\usepackage{iftex}
\ifPDFTeX
  \usepackage[T1]{fontenc}
  \usepackage[utf8]{inputenc}
  \usepackage{textcomp} % provide euro and other symbols
\else % if luatex or xetex
  \usepackage{unicode-math} % this also loads fontspec
  \defaultfontfeatures{Scale=MatchLowercase}
  \defaultfontfeatures[\rmfamily]{Ligatures=TeX,Scale=1}
\fi
\usepackage{lmodern}
\ifPDFTeX\else
  % xetex/luatex font selection
  \setmainfont[]{Helvetica Neue}
  \setmonofont[]{Menlo}
\fi
% Use upquote if available, for straight quotes in verbatim environments
\IfFileExists{upquote.sty}{\usepackage{upquote}}{}
\IfFileExists{microtype.sty}{% use microtype if available
  \usepackage[]{microtype}
  \UseMicrotypeSet[protrusion]{basicmath} % disable protrusion for tt fonts
}{}
\makeatletter
\@ifundefined{KOMAClassName}{% if non-KOMA class
  \IfFileExists{parskip.sty}{%
    \usepackage{parskip}
  }{% else
    \setlength{\parindent}{0pt}
    \setlength{\parskip}{6pt plus 2pt minus 1pt}}
}{% if KOMA class
  \KOMAoptions{parskip=half}}
\makeatother
\usepackage{longtable,booktabs,array}
\newcounter{none} % for unnumbered tables
\usepackage{calc} % for calculating minipage widths
% Correct order of tables after \paragraph or \subparagraph
\usepackage{etoolbox}
\makeatletter
\patchcmd\longtable{\par}{\if@noskipsec\mbox{}\fi\par}{}{}
\makeatother
% Allow footnotes in longtable head/foot
\IfFileExists{footnotehyper.sty}{\usepackage{footnotehyper}}{\usepackage{footnote}}
\makesavenoteenv{longtable}
\setlength{\emergencystretch}{3em} % prevent overfull lines
\providecommand{\tightlist}{%
  \setlength{\itemsep}{0pt}\setlength{\parskip}{0pt}}
% Header includes for docs/erw-model.tex (injected by build-erw-model.sh).
% Tighten longtable layout so wide tables with monospace cells fit on A4.
\usepackage{etoolbox}
\setlength{\tabcolsep}{4pt}
\renewcommand{\arraystretch}{1.05}
% Allow line breaks inside long \texttt{} strings (filenames, paths)
\usepackage{seqsplit}
% Inside longtable cells: shrink font and let \texttt{} break anywhere so long
% filenames (e.g. cdr_{climate_factor,pedogenic_fraction,per_t_basalt}.tif)
% wrap inside narrow columns instead of overflowing the page. seqsplit only
% inserts breakpoints — no visual effect unless a break is actually needed.
\protected\def\erwBreakableTexttInTables{%
  \let\origtexttt\texttt
  \def\texttt##1{\origtexttt{\seqsplit{##1}}}%
}
\AtBeginEnvironment{longtable}{\footnotesize\erwBreakableTexttInTables}
% Give the paragraph builder room to stretch before resorting to overfull lines.
\setlength{\emergencystretch}{3em}
\usepackage{bookmark}
\IfFileExists{xurl.sty}{\usepackage{xurl}}{} % add URL line breaks if available
\urlstyle{same}
\hypersetup{
  colorlinks=true,
  linkcolor={NavyBlue},
  filecolor={Maroon},
  citecolor={Blue},
  urlcolor={NavyBlue},
  pdfcreator={LaTeX via pandoc}}

\title{Carbon Dioxide Removal --- Technical Derivation}
\usepackage{etoolbox}
\makeatletter
\providecommand{\subtitle}[1]{% add subtitle to \maketitle
  \apptocmd{\@title}{\par {\large #1 \par}}{}{}
}
\makeatother
\subtitle{Per-pixel CDR calculation in the ex-ante ERW pipeline, with sourced mechanisms and constants}
\author{}
\date{2026-07-08}

\begin{document}
\maketitle

{
\hypersetup{linkcolor=NavyBlue}
\setcounter{tocdepth}{3}
\tableofcontents
}
\section{Carbon Dioxide Removal in the Ex-Ante ERW
Pipeline}\label{carbon-dioxide-removal-in-the-ex-ante-erw-pipeline}

\subsection{A technical derivation of the per-pixel CDR calculation,
with sources for every mechanism and
constant}\label{a-technical-derivation-of-the-per-pixel-cdr-calculation-with-sources-for-every-mechanism-and-constant}

This report documents exactly how the pipeline turns a tonne of applied
basalt into a quantity of durably removed carbon dioxide. It walks the
full chain implemented in \texttt{erw/erw-3-basalt-requirements.R}
(feedstock chemistry and the per-tonne potential) and
\texttt{erw/erw-9-cdr-yield.R} (the spatial climate, soil, and
accounting terms), states each governing equation, gives the numerical
default, and cites the peer-reviewed source that grounds each assumption
and mechanism.

Every bracketed reference \texttt{{[}n{]}} resolves in the
\textbf{References} section. Where a value is a modelling choice rather
than a measured constant, this is stated plainly and the choice is
defended against the cited range.

\begin{center}\rule{0.5\linewidth}{0.5pt}\end{center}

\subsection{1. Overview --- the calculation as one
chain}\label{overview-the-calculation-as-one-chain}

Durable CDR (tonnes CO₂ per hectare) at a cropland pixel is built in
five conceptual moves, then corrected by a sixth:

\begin{enumerate}
\def\labelenumi{\arabic{enumi}.}
\tightlist
\item
  \textbf{Stoichiometric ceiling} --- the maximum CO₂ a tonne of basalt
  could remove, fixed by its Ca and Mg oxide content (pure chemistry).
\item
  \textbf{Reactive fraction} --- only a fraction of that ceiling
  actually weathers on a decadal horizon.
\item
  \textbf{Grain-size factor} --- finer grinding exposes more surface
  area and speeds reaction.
\item
  \textbf{Climate \& soil-pH factor} --- temperature, moisture, and soil
  acidity rescale the reaction rate relative to a temperate reference
  site.
\item
  \textbf{Pedogenic-carbonate deduction} --- in dry soils some
  alkalinity re-precipitates as carbonate locally, re-releasing CO₂.
\item
  \textbf{Net-export correction} --- alkalinity consumed neutralizing
  the soil's own exchangeable acidity re-releases its CO₂ (as
  agricultural lime does) and is \emph{not} durable removal; only
  alkalinity exported beyond the acidity demand counts.
\end{enumerate}

Assembled, the per-hectare durable removal is

\[
\mathrm{CDR}_{\text{net}}(x)
= \max\!\Big(\underbrace{D(x)\cdot \tfrac{1}{1000}\,\Phi(x)\,C_{\max}\,(1-p_{\text{ped}}(x))}_{\text{gross removal}}
\;-\; \underbrace{F\cdot S(x)}_{\text{acidity re-release}}\;,\; 0\Big)
\]

where \(D\) is the basalt application rate (t ha⁻¹), \(C_{\max}\) the
stoichiometric ceiling (kg CO₂ t⁻¹), \(\Phi\) the effective reactive
fraction (grain × climate × pH), \(p_{\text{ped}}\) the
pedogenic-carbonate fraction, \(F\) the CO₂ re-released per tonne of
acidity neutralized, and \(S\) the soil's acidity expressed as a lime
requirement (t CaCO₃-eq ha⁻¹). Each term is derived below.

\subsubsection{Symbols and default
values}\label{symbols-and-default-values}

{\def\LTcaptype{none} % do not increment counter
\begin{longtable}[]{@{}
  >{\raggedright\arraybackslash}p{(\linewidth - 8\tabcolsep) * \real{0.2000}}
  >{\raggedright\arraybackslash}p{(\linewidth - 8\tabcolsep) * \real{0.2000}}
  >{\raggedright\arraybackslash}p{(\linewidth - 8\tabcolsep) * \real{0.2000}}
  >{\raggedright\arraybackslash}p{(\linewidth - 8\tabcolsep) * \real{0.2000}}
  >{\raggedright\arraybackslash}p{(\linewidth - 8\tabcolsep) * \real{0.2000}}@{}}
\toprule\noalign{}
\begin{minipage}[b]{\linewidth}\raggedright
Symbol
\end{minipage} & \begin{minipage}[b]{\linewidth}\raggedright
Meaning
\end{minipage} & \begin{minipage}[b]{\linewidth}\raggedright
Default
\end{minipage} & \begin{minipage}[b]{\linewidth}\raggedright
Units
\end{minipage} & \begin{minipage}[b]{\linewidth}\raggedright
Set in
\end{minipage} \\
\midrule\noalign{}
\endhead
\bottomrule\noalign{}
\endlastfoot
\(f_{\text{CaO}}, f_{\text{MgO}}\) & basalt oxide mass fractions & 0.10
/ 0.07 & -- & \texttt{erw-3} \\
\(C_{\max}\) & stoichiometric CDR ceiling & 309.8 & kg CO₂ / t basalt &
\texttt{erw-3} \\
\(\phi_0\) & reference reactive fraction & 0.20 & -- & \texttt{erw-3} \\
\(d\) & grain size & 50 & µm & \texttt{erw-3} \\
\(g\) & grain-size factor & 1.414 & -- & \texttt{erw-3} \\
\(\phi_{\text{ref}}=\phi_0 g\) & reference reactive fraction
(grain-adjusted) & 0.283 & -- & \texttt{erw-9} \\
\(T_{\text{ref}}, P_{\text{ref}}\) & climate reference anchors & 11 /
1000 & °C / mm & \texttt{erw-9} \\
\(\Phi\) & effective reactive fraction (capped at 1) & per pixel & -- &
\texttt{erw-9} \\
\(p_{\text{ped}}\) & pedogenic-carbonate fraction & 0--0.90 & -- &
\texttt{erw-9} \\
\(F\) & CO₂ re-released per t CaCO₃-eq neutralized & 0.88 & t CO₂ / t
CaCO₃ & \texttt{erw-9} \\
\(S\) & soil acidity as lime requirement & per pixel & t CaCO₃ / ha &
\texttt{erw-3} \\
\end{longtable}
}

\begin{center}\rule{0.5\linewidth}{0.5pt}\end{center}

\subsection{2. Feedstock chemistry --- the Ca and Mg
content}\label{feedstock-chemistry-the-ca-and-mg-content}

\textbf{What the model assumes.} A representative tropical-belt basalt
fine with \(f_{\text{CaO}}=0.10\) and \(f_{\text{MgO}}=0.07\) by mass
(\texttt{erw-3}, the \texttt{basalt} list). Minor oxides (K₂O 1.5 \%,
Na₂O 2.5 \%, P₂O₅ 0.3 \%) and trace metals (Ni 100 ppm, Cr 200 ppm) are
carried for nutrient and risk screening but do not enter the CDR term.

\textbf{Why these values.} Basalts carry their weathering ``fuel'' as
CaO ≈ 8--12 \% and MgO ≈ 5--10 \% by mass. For calibration, the six
commercial ERW basalts assayed by Lewis et al.~(2021) {[}1{]} (their
Table 6, p.~10) report \emph{elemental} Ca = 4.08--6.72 \% and Mg =
1.30--5.86 \%, i.e.~CaO ≈ 5.7--9.4 \% and MgO ≈ 2.2--9.7 \% --- so the
model's 10 \% CaO sits at/just above the top of that particular set and
7 \% MgO mid-range, consistent with the broader basalt population
(Renforth 2012 {[}44{]}, Abstract p.~229; compiled GEOROC dataset in
Lewis 2021 {[}1{]}). The pair is a representative default, exposed as an
override once quarry-specific assays are available.

\begin{center}\rule{0.5\linewidth}{0.5pt}\end{center}

\subsection{3. The stoichiometric CDR
ceiling}\label{the-stoichiometric-cdr-ceiling}

\textbf{Mechanism.} When a Ca- or Mg-silicate weathers under a
CO₂-charged soil solution, each mole of divalent base cation released is
charge-balanced by two moles of bicarbonate:

\[
\text{CaSiO}_3 + 2\,\text{CO}_2 + 3\,\text{H}_2\text{O}
\;\rightarrow\;
\text{Ca}^{2+} + 2\,\text{HCO}_3^- + \text{H}_4\text{SiO}_4 .
\]

So the ceiling is \textbf{2 mol CO₂ per mol of Ca²⁺ or Mg²⁺}. Using
molar masses (CaO 56.08, MgO 40.30, CO₂ 44.01 g mol⁻¹):

\[
\text{1 g CaO} \rightarrow \frac{2\times 44.01}{56.08} = 1.570\ \text{g CO}_2,
\qquad
\text{1 g MgO} \rightarrow \frac{2\times 44.01}{40.30} = 2.185\ \text{g CO}_2 .
\]

Per tonne of basalt,

\[
C_{\max} = \big(f_{\text{CaO}}\cdot 1.570 + f_{\text{MgO}}\cdot 2.185\big)\times 1000
= (0.10\cdot1.570 + 0.07\cdot2.185)\times1000 \approx \mathbf{309.8\ kg\ CO_2\,t^{-1}} .
\]

\textbf{Source.} The 1:2 cation-to-CO₂ bicarbonate stoichiometry is the
mechanistic basis of ERW CDR accounting (Hartmann et al.~2013 {[}45{]}).
The cation-based ``maximum uptake'' calculation from oxide chemistry,
and the resulting \textbf{≈ 0.3 t CO₂ per t basalt} for basic rocks,
trace to Renforth (2012) {[}44{]} and are used at scale by Beerling et
al.~(2020) {[}2{]} and Strefler et al.~(2018) {[}4{]}. It is an
\textbf{upper bound} in two senses: it assumes complete dissolution, and
it credits the full 2 mol CO₂ per cation, which only holds while carbon
is carried as dissolved bicarbonate --- on downstream carbonate
precipitation (Ca²⁺ + 2 HCO₃⁻ \(\rightarrow\) CaCO₃ + CO₂ + H₂O) half is
re-released, so the long-term net can fall toward \textasciitilde1 mol
CO₂ per mol Ca (Renforth 2019 {[}15{]}). Both reductions are applied
below (pedogenic deduction, §7; net-export, §9).

\begin{center}\rule{0.5\linewidth}{0.5pt}\end{center}

\subsection{4. The reactive (effective)
fraction}\label{the-reactive-effective-fraction}

\textbf{What the model assumes.} Only \(\phi_0 = 0.20\) of the
stoichiometric ceiling is realized on the accreditation horizon at the
reference grain size and climate (\texttt{erw-3},
\texttt{reactive\_fraction\_ref}).

\textbf{Why.} Field and mesocosm studies consistently report realized
CDR far below the stoichiometric ceiling because dissolution is
incomplete on decadal timescales, secondary minerals sequester cations,
and cation exchange delays export. Lewis et al.~(2021) {[}1{]} report
CDR of 1.3--8.5 t CO₂ ha⁻¹ over 15 yr at 50 t ha⁻¹ ---
i.e.~\textbf{26--170 kg CO₂ per tonne basalt actually realized}, against
the \textasciitilde310 kg ceiling. Their per-basalt spread (Cragmill ≈ 9
\%, Blue Ridge ≈ 21 \%, Tichum ≈ 57 \% of stoichiometric over 15 yr)
brackets \(\phi_0=0.20\) as a central value. Basalt field trials
(Kantola et al.~2023 {[}7{]}; Larkin et al.~2022 {[}8{]}) and the
uncertainty analyses of Reershemius \& Suhrhoff (2024) {[}16{]}) sit
within this band. A cool-temperate reality check sits at the low end:
Dupla et al.~(2025) {[}46{]} found detectable dissolution but
CO₂-removal rates \textless{} 10 \% of some other studies over
\textasciitilde1000 days of Swiss vineyard trials. The model therefore
treats \(\phi_0=0.20\) as a defensible central anchor (not a single
measured constant), later scaled up in warm/wet climates via the climate
factor and swept in \texttt{erw-8}.

\begin{center}\rule{0.5\linewidth}{0.5pt}\end{center}

\subsection{5. The grain-size factor}\label{the-grain-size-factor}

\textbf{Mechanism.} Dissolution is surface-reaction-limited, so rate
scales with specific surface area, which for a given mass scales
inversely with particle diameter. The model uses

\[
g = \left(\frac{d_{\text{ref}}}{d}\right)^{\beta},
\qquad d_{\text{ref}}=100\ \mu\text{m},\ \ \beta=0.5,
\]

giving \(g=(100/50)^{0.5}=\mathbf{1.414}\) at the 50 µm default.

\textbf{Why \(\beta=0.5\).} Empirically, silt-dominated basalt fines
(\textless45 µm) weather at roughly twice the rate of sand-dominated
material (150--500 µm), implying \(\beta\approx0.35\)--\(0.5\). The
square-root value is the parameter-free midpoint (rate ∝ √surface-area)
and avoids introducing an extra free constant. Grain-size control on
weathering rate and the associated grinding-energy penalty are
quantified by Strefler et al.~(2018) {[}4{]} (E ≈ 0.07 GJ t⁻¹ at 50 µm
rising to ≈ 3 GJ t⁻¹ at 2 µm) and used in the ERW potential estimates of
Beerling et al. (2020) {[}2{]}. Grain size is the model's dominant
decision variable and is swept 10--500 µm in \texttt{erw-8}.

\textbf{Reference reactive fraction after grinding.} \texttt{erw-9}
combines the two:

\[
\phi_{\text{ref}} = \phi_0\, g = 0.20\times 1.414 = \mathbf{0.283}.
\]

At the reference climate this yields an effective per-tonne removal of
\(\phi_{\text{ref}}C_{\max}=0.283\times309.8\approx \mathbf{87.6\ kg\ CO_2\,t^{-1}}\)
(\texttt{CDR\_eff\_kg\_per\_t\_ref}), consistent with the Lewis (2021)
mesocosm midpoint.

\begin{center}\rule{0.5\linewidth}{0.5pt}\end{center}

\subsection{6. The climate and soil-pH
factor}\label{the-climate-and-soil-ph-factor}

The reference reactive fraction is anchored to a temperate-humid site
(US Corn Belt, \(T_{\text{ref}}=11\) °C, \(P_{\text{ref}}=1000\) mm). A
dimensionless factor rescales it per pixel:

\[
\Phi(x) = \min\!\Big(\phi_{\text{ref}}\cdot f_{\text{MAT}}\cdot f_{\text{MAP}}\cdot f_{\text{pH}},\ 1\Big),
\]

\[
f_{\text{MAT}} = \mathrm{clamp}\!\left(\tfrac{\text{MAT}}{11},\,0.3,\,3.0\right),
\qquad
f_{\text{MAP}} = \mathrm{clamp}\!\left(\tfrac{\text{MAP}}{1000},\,0.3,\,3.0\right).
\]

\textbf{Temperature term.} Silicate dissolution is thermally activated
(Arrhenius), so warmer soils weather faster; the linear-in-\(T\)
multiplier is a first-order surrogate for that acceleration.
Basalt-specific dissolution experiments give the activation energies
directly: Gudbrandsson et al.~(2011) {[}40{]} measured crystalline
basalt over 5--75 °C, and Gíslason \& Oelkers (2003) {[}41{]} basaltic
glass over 6--150 °C, both confirming strong Arrhenius acceleration;
Palandri \& Kharaka (2004) {[}42{]} tabulate the activation energies
used in reactive-transport ERW models. The general temperature control
on weathering flux is established by White \& Brantley (2003) {[}11{]}
and Brantley et al.~(2008) {[}12{]}, and the climate control on ERW CDR
capacity specifically is mapped by Baek et al.~(2023) {[}5{]}.

\textbf{Moisture term.} Weathering export requires water moving through
the profile; CDR flux rises with precipitation and drainage. The
canonical field demonstration is White \& Blum (1995) {[}43{]}, who show
Si and Na weathering fluxes rising systematically with runoff across
watersheds; the dependence underpins the global ERW capacity maps of
Baek et al.~(2023) {[}5{]} and Beerling et al.~(2020) {[}2{]}.

\textbf{pH term (acid catalysis).} Silicate dissolution is
acid-catalyzed: the rate rises as pH falls (higher H⁺ activity), so
acidic soils dissolve basalt faster. The standard three-mechanism rate
law,
\(r = k_{\text{acid}}\,a_{\text{H}^+}^{\,n} + k_{\text{neutral}} + k_{\text{base}}\,a_{\text{H}^+}^{-m}\),
is compiled by Palandri \& Kharaka (2004) {[}42{]}; basalt-specific pH
data (Gíslason \& Oelkers 2003 {[}41{]}; Gudbrandsson et al.~2011
{[}40{]}) show dissolution rate rising sharply below neutral pH. The
model uses a triangular factor peaking near pH 5:

\[
f_{\text{pH}} =
\begin{cases}
0.5 & \text{pH} < 4\\
0.5 + 0.5(\text{pH}-4) & 4 \le \text{pH} < 5\\
1.0 - 0.25(\text{pH}-5) & 5 \le \text{pH} < 7\\
0.5 & \text{pH} \ge 7
\end{cases}
\]

The rise from pH 4\$\rightarrow\$5 and the \textasciitilde2× suppression
on calcareous soils (pH \textgreater{} 7) follow the acid-catalyzed
dissolution kinetics of White \& Brantley (2003) {[}11{]} and Brantley
et al.~(2008) {[}12{]}. The peak near pH 5 and the roll-off below it is
a modelling compromise: true rate is monotone in H⁺, but very low pH in
these soils co-occurs with Al toxicity and low base status, so the
factor is capped. This whole factor is flagged in code as a \textbf{v1
empirical proxy} to be replaced by a reactive-transport surrogate
(Kanzaki et al.~2023 {[}10{]}) when calibration data allow.

\begin{center}\rule{0.5\linewidth}{0.5pt}\end{center}

\subsection{7. The pedogenic-carbonate
deduction}\label{the-pedogenic-carbonate-deduction}

\textbf{Mechanism.} In semi-arid and arid soils a substantial fraction
of the Ca/Mg alkalinity released by weathering precipitates
\emph{locally} as secondary (pedogenic) CaCO₃/MgCO₃ instead of exporting
as bicarbonate. Carbonate precipitation releases roughly half the CO₂
the silicate originally consumed (Ca²⁺ + 2 HCO₃⁻ \(\rightarrow\) CaCO₃ +
CO₂ + H₂O), so this alkalinity does not deliver durable removal and must
be deducted:

\[
\text{exported fraction} = 1 - p_{\text{ped}}(x).
\]

\textbf{Keying on aridity.} \(p_{\text{ped}}\) is binned on the Aridity
Index (AI = MAP / PET) using UNEP (1992) classes {[}13{]}:

{\def\LTcaptype{none} % do not increment counter
\begin{longtable}[]{@{}lll@{}}
\toprule\noalign{}
AI = MAP/PET & Class & \(p_{\text{ped}}\) \\
\midrule\noalign{}
\endhead
\bottomrule\noalign{}
\endlastfoot
\textless{} 0.05 & Hyperarid & 0.90 \\
0.05--0.20 & Arid & 0.70 \\
0.20--0.50 & Semi-arid & 0.40 \\
0.50--0.65 & Dry sub-humid & 0.15 \\
\textgreater{} 0.65 & Humid & 0.00 \\
\end{longtable}
}

When the CGIAR-CSI Global Aridity Index raster (Zomer et al.~2022
{[}37{]}) is present the deduction is keyed directly on AI; otherwise
\texttt{erw-9} falls back to a MAP-only proxy with the same class
boundaries, which over-deducts in cool highlands (Ethiopia, Kenya,
Rwanda) where low PET keeps true AI humid.

\textbf{Sources.} Renforth (2019) {[}15{]} quantifies the ``≈ 0.5 mol
CO₂ re-released per mol CaCO₃ precipitated'' penalty; Harrington et
al.~(2023) {[}47{]} measure a concrete 16--27 \% CDR reduction from
carbonate precipitation during riverine transport; Kanzaki et al.~(2023)
{[}10{]} treat it as a permanence-loss term; and the Beerling et
al.~(2020) supplementary accounting {[}2{]} flags it. AI bands and
dataset: UNEP (1992) {[}13{]}; Zomer et al.~(2022) {[}37{]}.

\textbf{Open question.} Whether ERW-formed pedogenic carbonate is
durable (\textasciitilde10⁴ yr) or transient is contested:
fertilizer-derived strong acids can re-dissolve it and re-release the
CO₂, in which case it should not be credited as storage at all. The
exact loss fraction and its aridity scaling are among the
least-constrained numbers in the whole calculation --- the aridity bins
here are a deliberately coarse v1 proxy, and reviewers will (rightly)
push on the loss coefficients.

\begin{center}\rule{0.5\linewidth}{0.5pt}\end{center}

\subsection{8. Assembling the per-tonne and per-hectare gross
removal}\label{assembling-the-per-tonne-and-per-hectare-gross-removal}

Per tonne of basalt at pixel \(x\):

\[
c(x) = \Phi(x)\, C_{\max}\,\big(1-p_{\text{ped}}(x)\big) \quad[\text{kg CO}_2\,\text{t}^{-1}],
\]

and the \textbf{gross} per-hectare removal for an application rate
\(D(x)\) (t ha⁻¹):

\[
\mathrm{CDR}_{\text{gross}}(x) = D(x)\cdot \frac{c(x)}{1000}\quad[\text{t CO}_2\,\text{ha}^{-1}].
\]

\(D(x)\) comes from \texttt{erw-3}: either the agronomically
\textbf{targeted} lime-equivalent rate (basalt-to-lime ratio
\(1/\text{effective\_NV}=1/0.0996\approx 10.04\) t basalt per t
CaCO₃-eq, applied to the LiTAS lime requirement) or a \textbf{uniform}
10/20/50 t ha⁻¹ counterfactual. The targeted rate uses the LiTAS
lime-requirement method of Aramburu Merlos et al.~(2023) {[}35{]};
uniform rates match the modelling convention of Beerling et al.~(2020)
{[}2{]}, Baek et al.~(2023) {[}5{]}, and Kantola et al.~(2023) {[}7{]}.

\begin{center}\rule{0.5\linewidth}{0.5pt}\end{center}

\subsection{9. The net-export correction --- the load-bearing accounting
step}\label{the-net-export-correction-the-load-bearing-accounting-step}

This is the single most consequential step and the main methodological
update in this version of the pipeline.

\textbf{The problem.} The gross removal above credits carbon for
\emph{all} weathered alkalinity. But on an acid soil, alkalinity has two
competing fates drawing on one pool:

\begin{itemize}
\tightlist
\item
  \textbf{Neutralize exchangeable acidity} (displace Al³⁺/H⁺, raise pH).
  This is the agronomic benefit --- but neutralizing acidity consumes
  the bicarbonate and \textbf{releases the CO₂ back to the atmosphere}
  (HCO₃⁻ + H⁺ \(\rightarrow\) H₂O + CO₂, with the protons supplied by
  Al³⁺ hydrolysis). No net carbon is removed. This is exactly why
  agricultural lime is \emph{not} a CDR technology even though it raises
  yield the same way.
\item
  \textbf{Export as bicarbonate} to rivers and the ocean. This is
  durable removal (\textgreater10⁴ yr).
\end{itemize}

Counting the acidity-consumed tranche as CDR double-counts the same
alkalinity --- once as yield benefit, once as removal.

\textbf{The fix.} Deduct the acidity sink, a mass balance:

\[
\boxed{\ \mathrm{CDR}_{\text{net}}(x) = \max\!\big(\mathrm{CDR}_{\text{gross}}(x) - F\cdot S(x),\ 0\big)\ }
\]

\begin{itemize}
\tightlist
\item
  \(S(x)\) = the soil's acidity as a lime requirement (t CaCO₃-eq ha⁻¹),
  already computed by \texttt{erw-3} (Kamprath / LiTAS).
  \textbf{Regime-dependent:} year-1/NPV use the \emph{standing}
  exchangeable acidity (large); equilibrium uses only the \emph{annual
  re-acidification} increment (small).
\item
  \(F\) = CO₂ re-released per tonne of CaCO₃-equivalent acidity
  neutralized. It is \textbf{not a tuned parameter} --- it falls out of
  the feedstock's own constants:
\end{itemize}

\[
F = \frac{\text{CDR\_eff}}{\text{effective\_NV}\times 1000}
  = \frac{87.6}{99.6} \approx \mathbf{0.88}\ \text{t CO}_2\,(\text{t CaCO}_3)^{-1}.
\]

\textbf{Consequences.} Because standing acidity is large, the year-1
deduction removes \textasciitilde80 \% of credited CDR; at equilibrium
only the small annual re-acidification is deducted, so most removal
survives. Durable CDR is therefore largely a \textbf{steady-state
(maintenance)} phenomenon, and year-1 ERW on acid soils is, to first
order, \emph{liming with a small carbon tail}.

\textbf{Sources.} The carbon-fate rule --- bicarbonate produced against
\emph{strong/mineral} acid (including exchangeable soil acidity) is a
net CO₂ source, only bicarbonate produced against \emph{carbonic} acid
is a sink --- is established in the agricultural-liming
carbon-accounting literature (West \& McBride 2005 {[}38{]}; Hamilton et
al.~2007 {[}39{]}, \emph{Global Biogeochemical Cycles}). The most direct
ERW statement of the mechanism is Dietzen \& Rosing (2023) {[}48{]},
whose protocol states that non-carbonic acids must be accounted for
whenever soil pH (H₂O) is below \textbf{6.3}: as pH falls, mineral
dissolution rises but CO₂-capture \emph{efficiency} falls because
H⁺/exchangeable acidity consumes the released alkalinity, so naïve
cation-based accounting over-credits CDR on acid soils. That only
\emph{exported} alkalinity is durable removal --- and that
cation-release mass balance is an upper bound needing a downstream-loss
correction --- is the central message of the ERW MRV literature
(Reershemius et al.~2023 {[}50{]}; Suhrhoff et al.~2024 {[}49{]}; Zhang
et al.~2022 {[}51{]}; Kanzaki et al.~2023 {[}10{]}; Levy et al.~2024
{[}32{]}).

\textbf{One caveat, stated plainly.} This is a first-order
\emph{sequential} bound: it assumes the acidity sink is filled before
any alkalinity exports. In reality weathering, neutralization, and
leaching proceed simultaneously, so some alkalinity exports even on acid
soils, and re-acidification continually regenerates the sink. The truth
lies between the gross (``all exports'') and net (``sink filled first'')
bounds; the pipeline writes both \texttt{cdr\_yield\_*} (gross) and
\texttt{cdr\_yield\_*\_netexport} (net) rasters and leads the carbon
case on the equilibrium regime, where the two converge.

\begin{center}\rule{0.5\linewidth}{0.5pt}\end{center}

\subsection{10. Regime dependence and time-resolved
crediting}\label{regime-dependence-and-time-resolved-crediting}

The same per-hectare removal is credited under three temporal regimes
(in \texttt{erw-7}):

\begin{itemize}
\tightlist
\item
  \textbf{Year-1} --- one up-front application; CDR accrues over a
  5-year phasing (30/25/20/15/10 \%, \textasciitilde80 \% by year 5),
  first-order-kinetics-shaped (Kanzaki et al.~2023 {[}10{]}; Beerling et
  al.~2020 SI {[}2{]}).
\item
  \textbf{NPV @ 10 \%} --- the 5-year CDR tail is discounted by
  \(\text{CDR\_NPV\_FACTOR}=\sum_t f_t/(1+r)^t \big/ \sum_t f_t \approx 0.79\).
\item
  \textbf{Equilibrium} --- steady-state maintenance dose only; the
  net-export sink is the small annual re-acidification, so durable
  removal is largest here.
\end{itemize}

\begin{center}\rule{0.5\linewidth}{0.5pt}\end{center}

\subsection{11. Worked example (illustrative
pixel)}\label{worked-example-illustrative-pixel}

A warm, wet, moderately acid cropland pixel: MAT = 25 °C, MAP = 1200 mm,
soil pH = 6.0, humid (\(p_{\text{ped}}=0\)), targeted rate \(D=15\) t
ha⁻¹, standing acidity \(S_{\text{std}}=3.0\) t CaCO₃ ha⁻¹, maintenance
\(S_{\text{maint}}=0.30\) t CaCO₃ ha⁻¹.

{\def\LTcaptype{none} % do not increment counter
\begin{longtable}[]{@{}lll@{}}
\toprule\noalign{}
Step & Computation & Result \\
\midrule\noalign{}
\endhead
\bottomrule\noalign{}
\endlastfoot
\(f_{\text{MAT}}\) & clamp(25/11) & 2.27 \\
\(f_{\text{MAP}}\) & clamp(1200/1000) & 1.20 \\
\(f_{\text{pH}}\) & \(1-0.25(6-5)\) & 0.75 \\
\(\Phi\) & min(0.283·2.27·1.20·0.75, 1) & 0.578 \\
\(c\) (per t) & 0.578·309.8·(1−0) & 179 kg t⁻¹ \\
Gross CDR & 15·179/1000 & 2.69 t ha⁻¹ \\
Net (year-1) & max(2.69 − 0.88·3.0, 0) & \textbf{0.05 t ha⁻¹} \\
Net (equilibrium) & max(2.69 − 0.88·0.30, 0) & \textbf{2.43 t ha⁻¹} \\
\end{longtable}
}

The collapse from 2.69 \(\rightarrow\) 0.05 t ha⁻¹ in year-1, and its
survival at equilibrium, is the net-export effect in miniature.

\begin{center}\rule{0.5\linewidth}{0.5pt}\end{center}

\subsection{12. Assumptions and
limitations}\label{assumptions-and-limitations}

{\def\LTcaptype{none} % do not increment counter
\begin{longtable}[]{@{}
  >{\raggedright\arraybackslash}p{(\linewidth - 6\tabcolsep) * \real{0.2500}}
  >{\raggedright\arraybackslash}p{(\linewidth - 6\tabcolsep) * \real{0.2500}}
  >{\raggedright\arraybackslash}p{(\linewidth - 6\tabcolsep) * \real{0.2500}}
  >{\raggedright\arraybackslash}p{(\linewidth - 6\tabcolsep) * \real{0.2500}}@{}}
\toprule\noalign{}
\begin{minipage}[b]{\linewidth}\raggedright
\#
\end{minipage} & \begin{minipage}[b]{\linewidth}\raggedright
Assumption / mechanism
\end{minipage} & \begin{minipage}[b]{\linewidth}\raggedright
Status
\end{minipage} & \begin{minipage}[b]{\linewidth}\raggedright
Upgrade path
\end{minipage} \\
\midrule\noalign{}
\endhead
\bottomrule\noalign{}
\endlastfoot
1 & 10 \%/7 \% CaO/MgO basalt & Representative; override per quarry &
Quarry assays via \texttt{erw-basalt-access} \\
2 & 1:2 cation:CO₂ stoichiometry & Well-established chemistry & --- \\
3 & 0.20 reference reactive fraction & Central of a wide (0.08--0.55)
field range & Site calibration as trials accrue \\
4 & \(\beta=0.5\) grain-size exponent & Defensible midpoint (0.35--0.5)
& Feedstock-specific SSA-rate data \\
5 & Linear climate multipliers & v1 empirical proxy & Reactive-transport
surrogate {[}10{]} \\
6 & Triangular pH factor & v1 proxy; true rate monotone in H⁺ & Kinetic
dissolution model \\
7 & Aridity-binned pedogenic loss & Coarse; MAP proxy over-deducts
highlands & CGIAR-CSI AI raster {[}37{]} \\
8 & Net-export sequential bound & First-order; presented with gross
bound & Coupled transport-reaction \\
\end{longtable}
}

The two quantities the downstream economics is most sensitive to --- the
realized reactive fraction (Step 4) and the pedogenic/net-export loss
(Steps 7 and 9) --- are also the least empirically resolved, which is
why they are carried explicitly as bounds and swept in \texttt{erw-8}
rather than fixed silently.

\begin{center}\rule{0.5\linewidth}{0.5pt}\end{center}

\subsection{13. Provenance of every number and
formula}\label{provenance-of-every-number-and-formula}

Each quantity in this report traces to one of three kinds of source: a
\textbf{standard physical constant}, an \textbf{in-model derivation or
modelling choice} (whose provenance is a specific line in the pipeline
code, not a journal page), or a \textbf{literature-anchored value}
(traced to an exact page / table / figure / equation, with a confidence
flag from the source-verification pass). Line numbers refer to the
scripts as committed. Confidence: \textbf{V} = verified from full text,
\textbf{A} = confirmed from the paper's abstract only (full text
paywalled), \textbf{C} = modelling choice / convention rather than a
single measured value.

\subsubsection{13.1 Standard physical
constants}\label{standard-physical-constants}

{\def\LTcaptype{none} % do not increment counter
\begin{longtable}[]{@{}
  >{\raggedright\arraybackslash}p{(\linewidth - 4\tabcolsep) * \real{0.3333}}
  >{\raggedright\arraybackslash}p{(\linewidth - 4\tabcolsep) * \real{0.3333}}
  >{\raggedright\arraybackslash}p{(\linewidth - 4\tabcolsep) * \real{0.3333}}@{}}
\toprule\noalign{}
\begin{minipage}[b]{\linewidth}\raggedright
Quantity
\end{minipage} & \begin{minipage}[b]{\linewidth}\raggedright
Value
\end{minipage} & \begin{minipage}[b]{\linewidth}\raggedright
Source and locus
\end{minipage} \\
\midrule\noalign{}
\endhead
\bottomrule\noalign{}
\endlastfoot
Molar mass CaO / MgO & 56.08 / 40.30 g mol⁻¹ & IUPAC/CIAAW standard
atomic weights 2021 {[}52{]} (Ca 40.078, Mg 24.305, O 15.999) \\
Molar mass CO₂ / CaCO₃ & 44.01 / 100.09 g mol⁻¹ & IUPAC/CIAAW 2021
{[}52{]} (C 12.011) \\
1 mol Ca²⁺/Mg²⁺ \(\rightarrow\) 2 mol CO₂ (bicarbonate basis) &
stoichiometry & Hartmann et al.~2013 {[}45{]}, \textbf{Eq. (8), §2.1}
(forsterite reaction) --- \textbf{V}; Beerling et al.~2020 {[}2{]},
Methods ``CDR'', Eqs (4)--(5) \\
\end{longtable}
}

\subsubsection{13.2 In-model derivations and modelling choices
(provenance = code
line)}\label{in-model-derivations-and-modelling-choices-provenance-code-line}

{\def\LTcaptype{none} % do not increment counter
\begin{longtable}[]{@{}
  >{\raggedright\arraybackslash}p{(\linewidth - 4\tabcolsep) * \real{0.3333}}
  >{\raggedright\arraybackslash}p{(\linewidth - 4\tabcolsep) * \real{0.3333}}
  >{\raggedright\arraybackslash}p{(\linewidth - 4\tabcolsep) * \real{0.3333}}@{}}
\toprule\noalign{}
\begin{minipage}[b]{\linewidth}\raggedright
Quantity
\end{minipage} & \begin{minipage}[b]{\linewidth}\raggedright
Value
\end{minipage} & \begin{minipage}[b]{\linewidth}\raggedright
Code locus
\end{minipage} \\
\midrule\noalign{}
\endhead
\bottomrule\noalign{}
\endlastfoot
CaO / MgO mass fraction & 0.10 / 0.07 & \texttt{erw-3.R:55–56} (choice;
§13.3) \\
Grain size \(d\) & 50 µm & \texttt{erw-3.R:67} (decision variable) \\
Reference reactive fraction \(\phi_0\) & 0.20 & \texttt{erw-3.R:74}
(anchored; §13.3) \\
Theoretical CCE (1.783, 2.481 g/g) & 0.352 & \texttt{erw-3.R:82–84}
(molar-mass arithmetic) \\
Stoichiometric ceiling \(C_{\max}\) (1.570, 2.185 g/g) & 309.8 kg t⁻¹ &
\texttt{erw-3.R:87–91} \\
Grain-size factor \(g=(100/d)^{0.5}\) & 1.414 & \texttt{erw-3.R:100–102}
(\(\beta=0.5\) choice; §13.3) \\
Grinding curve \(E=0.07(50/d)^{1.2}\) & --- & \texttt{erw-3.R:113–118}
(anchored to Strefler; §13.3) \\
Effective NV & 0.0996 t t⁻¹ & \texttt{erw-3.R:130} \\
Basalt-to-lime ratio \(1/\text{NV}\) & 10.04 & \texttt{erw-3.R:131} \\
Reference effective CDR & 87.6 kg t⁻¹ & \texttt{erw-3.R:132} \\
\(\phi_{\text{ref}}=\phi_0 g\) & 0.283 & \texttt{erw-9.R:48} \\
\(T_{\text{ref}}, P_{\text{ref}}\) & 11 °C, 1000 mm &
\texttt{erw-9.R:83–84} (US Corn Belt normal) \\
\(f_{\text{MAT}}, f_{\text{MAP}}\) clamp {[}0.3, 3.0{]} & --- &
\texttt{erw-9.R:86–87} \\
\(f_{\text{pH}}\) triangular, peak ≈ 5 & --- & \texttt{erw-9.R:90–93} \\
\(\text{climate\_factor}=f_{\text{MAT}}f_{\text{MAP}}f_{\text{pH}}\) &
--- & \texttt{erw-9.R:96} \\
Pedogenic fractions 0.90/0.70/0.40/0.15/0 & --- &
\texttt{erw-9.R:152–163} (bands from §13.3) \\
Reactive fraction cap at 1 & --- & \texttt{erw-9.R:181} \\
Per-tonne CDR \(c(x)\) & --- & \texttt{erw-9.R:186} \\
CDR yield \(= D\cdot c/1000\) & --- & \texttt{erw-9.R:194} \\
\(F=\text{CDR\_eff}/(\text{NV}\times1000)\) & 0.88 &
\texttt{erw-9.R:233} \\
Net-export \(\max(\text{gross}-F S,0)\) & --- &
\texttt{erw-9.R:240–244} \\
CDR phasing 30/25/20/15/10 \% & --- & \texttt{erw-7.R:106} (shape from
§13.3) \\
Discount rate & 10 \% & \texttt{erw-7.R:100} \\
CDR NPV factor & ≈ 0.79 & \texttt{erw-7.R:109–111} \\
\end{longtable}
}

\subsubsection{13.3 Literature-anchored values (source, exact locus,
confidence)}\label{literature-anchored-values-source-exact-locus-confidence}

{\def\LTcaptype{none} % do not increment counter
\begin{longtable}[]{@{}
  >{\raggedright\arraybackslash}p{(\linewidth - 4\tabcolsep) * \real{0.3333}}
  >{\raggedright\arraybackslash}p{(\linewidth - 4\tabcolsep) * \real{0.3333}}
  >{\raggedright\arraybackslash}p{(\linewidth - 4\tabcolsep) * \real{0.3333}}@{}}
\toprule\noalign{}
\begin{minipage}[b]{\linewidth}\raggedright
Quantity
\end{minipage} & \begin{minipage}[b]{\linewidth}\raggedright
Source --- exact locus
\end{minipage} & \begin{minipage}[b]{\linewidth}\raggedright
Conf.
\end{minipage} \\
\midrule\noalign{}
\endhead
\bottomrule\noalign{}
\endlastfoot
Representative basalt CaO/MgO & Lewis 2021 {[}1{]}, \textbf{Table 6,
p.~10} (elemental Ca 4.08--6.72 \%, Mg 1.30--5.86 \%); Renforth 2012
{[}44{]}, \textbf{Abstract, p.~229} & V / A \\
Stoichiometric max ≈ 0.3 t CO₂ t⁻¹ (basic rock) & Renforth 2012
{[}44{]}, \textbf{Abstract, p.~229}; Lewis 2021 {[}1{]}, \textbf{Table 6
p.~10} (\(R_{\text{CO}_2}\)(CaO+MgO) = 0.13--0.30), formula \textbf{Eq.
(5), p.~13} & A / V \\
Reactive fraction \textasciitilde0.20 (realized/stoichiometric) & Lewis
2021 {[}1{]}, \textbf{Table 5, p.~5} (15-yr CDR 1.3--8.5 t ha⁻¹ at 50 t
ha⁻¹: Cragmill 1.3 \(\rightarrow\) Tichum 8.5; ≈ 9--57 \% of
stoichiometric); Abstract p.~1; §4.1 p.~13 & V \\
Low-end realized fraction (\textless{} 10 \%) & Dupla et al.~2025
{[}46{]}, \textbf{Abstract} & A \\
Grinding energy 0.07 GJ t⁻¹ @ 50 µm; 3 GJ t⁻¹ @ 2 µm & Strefler 2018
{[}4{]}, \textbf{``Mining, crushing and grinding'' subsection
(\textasciitilde p.~7); Fig. 1 (lower axis); SI C, Eq. C1} & V \\
Grain-size \(\rightarrow\) dissolution-rate scaling
(\(\beta\approx0.35\)--0.5) & Strefler 2018 {[}4{]}, \textbf{Eq. (3) \&
Fig. 1 (\textasciitilde p.~3); Eqs (1)--(2), SI B} (\(\beta\) is a
modelling convention grounded in surface-area kinetics) & V / C \\
Temperature (Arrhenius) dependence & Gudbrandsson et al.~2011 {[}40{]};
Gíslason \& Oelkers 2003 {[}41{]}; Palandri \& Kharaka 2004 {[}42{]}
(activation-energy tables) & V \\
Moisture / runoff dependence & White \& Blum 1995 {[}43{]} & V \\
pH acid-catalysis rate law & Palandri \& Kharaka 2004 {[}42{]}
(three-mechanism rate law); Gíslason \& Oelkers 2003 {[}41{]}; White \&
Brantley 2003 {[}11{]} & V \\
Aridity bands 0.05/0.20/0.50/0.65 & UNEP 1992 {[}13{]}, reproduced
verbatim in IPCC SRCCL 2019 {[}53{]}, \textbf{Ch. 3 §3.1.1, Fig. 3.1,
p.~\textasciitilde255} (note: Zomer 2022 {[}37{]} uses a 0.03 hyperarid
variant --- not used here) & V \\
Pedogenic carbonate re-releases ≈ 0.5 mol CO₂ per mol CaCO₃ & Renforth
2019 {[}15{]}, \textbf{Eq. (3), Introduction}; Beerling 2020 {[}2{]},
Methods ``CDR'', \textbf{Eq. (5)} (1 mol CO₂/mol Ca vs 1.72 via ocean) &
V \\
Carbonate-precipitation CDR loss 16--27 \% & Harrington et al.~2023
{[}47{]}, \textbf{Abstract} & A \\
Acidity-consumed alkalinity ≠ CDR; pH (H₂O) 6.3 threshold & Dietzen \&
Rosing 2023 {[}48{]}, \textbf{Abstract} & A \\
Ocean-transport permanence / leakage (\textasciitilde9 \% silicate) &
Kanzaki et al.~2023 {[}10{]}, \textbf{Results/Discussion, Fig. 2} & V \\
LiTAS lime-requirement method (targeted rate) & Aramburu Merlos 2023
{[}35{]}, \textbf{§4.1.6, Eqs (11)--(12)} (\(a=0.60,\ b=0.92\)) & V \\
Uniform rates 10/20/50 t ha⁻¹ & Beerling 2020 {[}2{]} baseline
\textbf{40 t ha⁻¹ yr⁻¹} (Methods ``Baseline simulations''); Baek 2023
{[}5{]} 10 t ha⁻¹; Kantola 2023 {[}7{]} 50 t ha⁻¹ & A \\
CDR phasing shape (first-order, \textasciitilde80 \% by yr 5) & Kanzaki
2023 {[}10{]}; Beerling 2020 {[}2{]} SI & C \\
\end{longtable}
}

\textbf{Note on two thresholds.} (i) The pipeline's hyperarid break is
0.05 (UNEP 1992 / IPCC SRCCL), \emph{not} the 0.03 in Zomer et
al.~(2022); if the CGIAR-CSI raster is loaded its own scaling is
respected but the class breaks stay at the code values. (ii) Renforth
(2012)'s ``\textasciitilde0.3 t CO₂ t⁻¹'' and its cation-uptake formula
could only be verified at the abstract; the in-text derivation is behind
a paywall and is therefore marked \textbf{A}.

\begin{center}\rule{0.5\linewidth}{0.5pt}\end{center}

\subsection{References}\label{references}

\begin{enumerate}
\def\labelenumi{\arabic{enumi}.}
\tightlist
\item
  Lewis, A. L., et al.~(2021). Effects of mineralogy, chemistry and
  physical properties of basalts on carbon capture potential and
  plant-nutrient element release via enhanced weathering. \emph{Applied
  Geochemistry}, 132, 105023.
  https://doi.org/10.1016/j.apgeochem.2021.105023
\item
  Beerling, D. J., et al.~(2020). Potential for large-scale CO₂ removal
  via enhanced rock weathering with croplands. \emph{Nature}, 583,
  242--248. https://doi.org/10.1038/s41586-020-2448-9
\item
  Strefler, J., Amann, T., Bauer, N., Kriegler, E., \& Hartmann, J.
  (2018). Potential and costs of carbon dioxide removal by enhanced
  weathering of rocks. \emph{Environmental Research Letters}, 13(3),
  034010. https://doi.org/10.1088/1748-9326/aaa9c4
\item
  Baek, S. H., et al.~(2023). Impact of climate on the global capacity
  for enhanced rock weathering on croplands. \emph{Earth's Future}, 11,
  e2023EF003698. https://doi.org/10.1029/2023EF003698
\item
  Kantola, I. B., et al.~(2023). Improved net carbon budgets in the U.S.
  Midwest through direct measured impacts of enhanced weathering.
  \emph{Global Change Biology}, 29(24), 7012--7028.
  https://doi.org/10.1111/gcb.16903
\item
  Larkin, C. S., et al.~(2022). Quantification of CO₂ removal in a
  large-scale enhanced weathering field trial on an oil palm plantation
  in Sabah, Malaysia. \emph{Frontiers in Climate}, 4, 959229.
  https://doi.org/10.3389/fclim.2022.959229
\item
  Kanzaki, Y., Planavsky, N. J., \& Reinhard, C. T. (2023). New
  estimates of the storage permanence and ocean co-benefits of enhanced
  rock weathering. \emph{PNAS Nexus}, 2(4), pgad059.
  https://doi.org/10.1093/pnasnexus/pgad059
\item
  White, A. F., \& Brantley, S. L. (2003). The effect of time on the
  weathering of silicate minerals. \emph{Chemical Geology}, 202(3--4),
  479--506. https://doi.org/10.1016/j.chemgeo.2003.03.001
\item
  Brantley, S. L., Kubicki, J. D., \& White, A. F. (Eds.) (2008).
  \emph{Kinetics of Water-Rock Interaction}. Springer.
  https://doi.org/10.1007/978-0-387-73563-4
\item
  UNEP (1992). \emph{World Atlas of Desertification}. United Nations
  Environment Programme, London: Edward Arnold.
\item
  Renforth, P. (2019). The negative emission potential of alkaline
  materials. \emph{Nature Communications}, 10, 1401.
  https://doi.org/10.1038/s41467-019-09475-5
\item
  Reershemius, T., \& Suhrhoff, T. J. (2024). On error, uncertainty, and
  assumptions in calculating carbon dioxide removal rates by enhanced
  rock weathering. \emph{Global Change Biology}, 30(1), e17025.
  https://doi.org/10.1111/gcb.17025
\item
  Levy, C. R., et al.~(2024). Enhanced rock weathering for carbon
  removal --- monitoring and mitigating potential environmental impacts
  on agricultural land. \emph{Environmental Science \& Technology},
  58(39), 17215--17226. https://doi.org/10.1021/acs.est.4c02368
\item
  Aramburu Merlos, F., Silva, J. V., Baudron, F., \& Hijmans, R. J.
  (2023). Estimating lime requirements for tropical soils: Model
  comparison and development. \emph{Geoderma}, 432, 116421.
  https://doi.org/10.1016/j.geoderma.2023.116421
\item
  Zomer, R. J., Xu, J., \& Trabucco, A. (2022). Version 3 of the Global
  Aridity Index and Potential Evapotranspiration Database.
  \emph{Scientific Data}, 9, 409.
  https://doi.org/10.1038/s41597-022-01493-1
\item
  West, T. O., \& McBride, A. C. (2005). The contribution of
  agricultural lime to CO₂ emissions in the United States: dissolution,
  transport, and net emissions. \emph{Agriculture, Ecosystems \&
  Environment}, 108(2), 145--154.
  https://doi.org/10.1016/j.agee.2005.01.002
\item
  Hamilton, S. K., Kurzman, A. L., Arango, C., Jin, L., \& Robertson, G.
  P. (2007). Evidence for carbon sequestration by agricultural liming.
  \emph{Global Biogeochemical Cycles}, 21, GB2021.
  https://doi.org/10.1029/2006GB002738
\item
  Gudbrandsson, S., Wolff-Boenisch, D., Gíslason, S. R., \& Oelkers, E.
  H. (2011). An experimental study of crystalline basalt dissolution
  from 2 ≤ pH ≤ 11 and temperatures from 5 to 75 °C. \emph{Geochimica et
  Cosmochimica Acta}, 75(19), 5496--5509.
  https://doi.org/10.1016/j.gca.2011.06.035
\item
  Gíslason, S. R., \& Oelkers, E. H. (2003). Mechanism, rates, and
  consequences of basaltic glass dissolution: II. An experimental study
  of the dissolution rates of basaltic glass as a function of pH and
  temperature. \emph{Geochimica et Cosmochimica Acta}, 67(20),
  3817--3832. https://doi.org/10.1016/S0016-7037(03)00176-5
\item
  Palandri, J. L., \& Kharaka, Y. K. (2004). \emph{A compilation of rate
  parameters of water--mineral interaction kinetics for application to
  geochemical modeling}. USGS Open-File Report 2004-1068.
  https://pubs.usgs.gov/publication/ofr20041068
\item
  White, A. F., \& Blum, A. E. (1995). Effects of climate on chemical
  weathering in watersheds. \emph{Geochimica et Cosmochimica Acta},
  59(9), 1729--1747. https://doi.org/10.1016/0016-7037(95)00078-E
\item
  Renforth, P. (2012). The potential of enhanced weathering in the UK.
  \emph{International Journal of Greenhouse Gas Control}, 10, 229--243.
  https://doi.org/10.1016/j.ijggc.2012.06.011
\item
  Hartmann, J., West, A. J., Renforth, P., Köhler, P., De La Rocha, C.
  L., Wolf-Gladrow, D. A., Dürr, H. H., \& Scheffran, J. (2013).
  Enhanced chemical weathering as a geoengineering strategy to reduce
  atmospheric carbon dioxide, supply nutrients, and mitigate ocean
  acidification. \emph{Reviews of Geophysics}, 51(2), 113--149.
  https://doi.org/10.1002/rog.20004
\item
  Dupla, X., Bertagni, M. B., \& Grand, S. (2025). Three years of field
  trials indicate a sustained enhanced rock weathering signal with
  limited CO₂ removal. \emph{Environmental Science \& Technology},
  59(48), 25751--25764. https://doi.org/10.1021/acs.est.5c09820
\item
  Harrington, K. J., Hilton, R. G., \& Henderson, G. M. (2023).
  Implications of the riverine response to enhanced weathering for CO₂
  removal in the UK. \emph{Applied Geochemistry}, 152, 105643.
  https://doi.org/10.1016/j.apgeochem.2023.105643
\item
  Dietzen, C., \& Rosing, M. T. (2023). Quantification of CO₂ uptake by
  enhanced weathering of silicate minerals applied to acidic soils.
  \emph{International Journal of Greenhouse Gas Control}, 125, 103872.
  https://doi.org/10.1016/j.ijggc.2023.103872
\item
  Suhrhoff, T. J., Reershemius, T., Wang, J., Jordan, J. S., Reinhard,
  C. T., \& Planavsky, N. J. (2024). Measuring enhanced weathering:
  inorganic carbon-based approaches may be required to complement
  cation-based approaches. \emph{Frontiers in Climate}, 6, 1352825.
  https://doi.org/10.3389/fclim.2024.1352825
\item
  Reershemius, T., Kelland, M. E., Jordan, J. S., et al.~(2023). Initial
  validation of a soil-based mass-balance approach for empirical
  monitoring of enhanced rock weathering rates. \emph{Environmental
  Science \& Technology}, 57(48), 19497--19507.
  https://doi.org/10.1021/acs.est.3c03609
\item
  Zhang, S., Planavsky, N. J., Katchinoff, J., Raymond, P. A., Kanzaki,
  Y., Reershemius, T., \& Reinhard, C. T. (2022). River chemistry
  constraints on the carbon capture potential of surficial enhanced rock
  weathering. \emph{Limnology \& Oceanography}, 67(S1), S148--S157.
  https://doi.org/10.1002/lno.12244
\item
  Prohaska, T., Irrgeher, J., Benefield, J., et al.~(2022). Standard
  atomic weights of the elements 2021 (IUPAC Technical Report).
  \emph{Pure and Applied Chemistry}, 94(5), 573--600.
  https://doi.org/10.1515/pac-2019-0603
\item
  IPCC (2019). Desertification (Chapter 3). In \emph{Climate Change and
  Land: an IPCC Special Report (SRCCL)}. Mirzabaev, A., Wu, J., Evans,
  J., et al.~https://www.ipcc.ch/srccl/chapter/chapter-3/
\end{enumerate}

\emph{References {[}1{]}--{[}37{]} match the master catalogue in
\texttt{docs/erw-parameters.md}. {[}38{]}--{[}53{]} were added and
DOI-verified for this report to source the weathering-kinetics,
stoichiometry, net-export, and constant-provenance items specifically.
Page/table/equation loci in §13.3 were located by direct inspection of
each source; items marked \textbf{A} were confirmable only at the
abstract because the full text is paywalled.}

\end{document}
